\clearpage
\newpage
\appendix
\section{Additional Implementation Details}
\label{sec:appendix_impl}

\cref{tab:hypers_all} summarizes the configurations and hyper-parameters for ablation studies and system-level comparisons. 
We provide detailed experimental configurations for reproducibility. All ablation studies share a common default setup, while system-level comparisons use scaled-up configurations. More implementation details are described as follows.

\input{tab/config}

\subsection{Pseudo-code for Computing Drifting Field $\V$}
\label{sec:appendix_compute_V}

Alg.~\ref{alg:drift_code} provides the pseudo-code for computing $\V$. The computation is based on taking empirical means in Eq.~(\ref{eq:VZZ}) and (\ref{eq:kernel}), which are implemented as softmax over $\y$-sample axis. In practice, we further normalize over the $\x$-sample axis, also implemented as softmax on the same logit matrix. 
We ablate its influence in \ref{sec:appendix_kernel_normalizatio}.

\renewcommand{\yn}{\y^{-}}
\renewcommand{\yp}{\y^{+}}
\newcommand{\B}{\mathcal{B}}

It is worth noting that this implementation preserves the desired property of $\V$.
In principle, this implementation can be viewed as a Monte Carlo estimation of a drifting field:
\begin{equation}
\V_{p,q}(\x)=\mathbb E_{\B,p,q}[\tilde K_\B(\x, \yp)\tilde K_\B(\x, \yn)(\yp-\yn)],    
\end{equation}
where $\B$ consists of other samples in the batch and $\tilde K_\B$ denote normalizing the distance based on statistics within $\B$. This $\V$ also satisfies $\V_{p,p}(\x)=\mathbf 0$, since when $p=q$, the term  $\tilde K_\B(\yp, x)\tilde K_\B(\yn,x)(\yp-\yn)$ cancels out with the term $\tilde K_\B(\yn, x)\tilde K_\B(\yp,x)(\yn-\yp)$. 

\begin{algorithm}[t]
\caption{Computing the drifting field $\V$.}
\label{alg:drift_code}
\begin{lstlisting}
def compute_V(x, y_pos, y_neg, T):
  # x: [N, D]
  # y_pos: [N_pos, D]
  # y_neg: [N_neg, D]
  # T: temperature

  # compute pairwise distance
  dist_pos = cdist(x, y_pos)  # [N, N_pos]
  dist_neg = cdist(x, y_neg)  # [N, N_neg]
  
  # ignore self (if y_neg is x)
  dist_neg += eye(N) * 1e6

  # compute logits
  logit_pos = -dist_pos / T
  logit_neg = -dist_neg / T

  # concat for normalization
  logit = cat([logit_pos, logit_neg], dim=1)
  # normalize along both dimensions
  A_row = logit.softmax(dim=-1)
  A_col = logit.softmax(dim=-2)
  A = sqrt(A_row * A_col) 

  # back to [N, N_pos] and [N, N_neg]
  A_pos, A_neg = split(A, [N_pos,], dim=1)

  # compute the weights
  W_pos = A_pos  # [N, N_pos]
  W_neg = A_neg  # [N, N_neg]
  W_pos *= A_neg.sum(dim=1,keepdim=True)
  W_neg *= A_pos.sum(dim=1,keepdim=True)

  drift_pos = W_pos @ y_pos # [N_x, D]
  drift_neg = W_neg @ y_neg # [N_x, D]

  V = drift_pos - drift_neg
  return V 
\end{lstlisting}
\end{algorithm}

\subsection{Generator Architecture}
\label{sec:appendix_generator}

\paragraph{Input and output.}
The input to the generator consists of random noise along with conditioning:
\[
f_\theta: (\e, c, \alpha) \mapsto \x
\]
where $\e$ denotes random variables, $c$ is a class label, and $\alpha$ is the CFG strength.
$\e$ may consist of both continuous random variables (\eg, Gaussian noise) and discrete ones (\eg, uniformly distributed integers; see random style embeddings).
For latent-space models, the output $\mathbf{x} \in \mathbb{R}^{32\times 32\times 4}$ is in the SD-VAE latent space.
For pixel-space models, the output $\mathbf{x} \in \mathbb{R}^{256\times 256\times 3}$ is directly an image.

\paragraph{Transformer.}
We adopt a DiT-style Transformer~\citep{peebles2023scalable}. Following~\citep{yao2025reconstruction}, we use SwiGLU~\citep{shazeer2020glu}, RoPE~\citep{su2024roformer}, RMSNorm~\citep{zhang2019root}, and QK-Norm~\citep{henry2020query}.
The input Gaussian noise is patchified into 256$=$16$\times$16 tokens (patch size 2$\times$2 for latent, 16$\times$16 for pixel).
Conditioning $(c, \alpha)$ is processed by adaLN, as well as by in-context conditioning tokens.
The output tokens are unpatchified back to the target shape.

\paragraph{In-context tokens.}
Following~\cite{li2025back}, we prepend 16 learnable tokens to the sequence for in-context conditioning \cite{peebles2023scalable}.
These tokens are formed by summing the projected conditioning vector with positional embeddings.
\paragraph{Random style embeddings.}
Our framework allows arbitrary noise distributions beyond Gaussians.
Inspired by StyleGAN~\cite{sauer2022stylegan}, we introduce an additional 32 ``style tokens'': each of which is a random index into a codebook of 64 learnable embeddings.
These are summed and added to the conditioning vector.
This does not change the sequence length and introduces negligible overhead in terms of parameters and FLOPs. This table reports the effect of style embeddings on our ablation default:

\begin{center}
\tablestyle{6pt}{1.02}
\tablefontsize
\begin{tabular}{c|cc}
\toprule
 & w/o style & w/ style \\
\midrule
FID & 8.86 & \textbf{8.46} \\
\bottomrule
\end{tabular}
\end{center}

In contrast to diffusion-/flow-based methods, our method can naturally handle different types of noise or random variables.
With random style embeddings, the input random variables consist of two parts: (1) Gaussian noise, and (2) discrete indices for style embeddings. Our model $f$ produces the pushforward distribution of their joint distribution. 

%============================================================
% MAE Implementation
%============================================================
\subsection{Implementation of ResNet-style MAE}
\label{sec:appendix_latent_mae}

In addition to standard self-supervised learning models (MoCo \cite{he2020momentum}, SimCLR\cite{chen2020simple}), we develop a customized ResNet-style MAE model as the feature encoder for drifting loss.

\paragraph{Overview.}
Unlike standard MAE \cite{he2022masked}, which is based on ViT \cite{Dosovitskiy2021}, our MAE trains a convolutional ResNet that provides multi-scale features.
For latent-space models, the input and output have dimension 32$\times$32$\times$4;
for pixel-space models, the input and output have dimension 256$\times$256$\times$3.

Our MAE consists of a ResNet-style encoder paired with a deconvolutional decoder in a U-Net-style \cite{ronneberger2015u} encoder-decoder architecture.
We only use the ResNet-style encoder for feature extraction when computing the drifting loss.

\paragraph{MAE Encoder.}
The encoder follows a classical ResNet \cite{he2016deep} design.
It maps an input to multi-scale feature maps (4 scales in ResNet):
\[
\text{Encoder}: \mathbf{x} \mapsto \{\mathbf{f}_1, \mathbf{f}_2, \mathbf{f}_3, \mathbf{f}_4\}
\]
Here, a feature map $\mathbf{f}_i$ has dimension $H_i{\times}W_i{\times}C_i$, with $H_i{\times}W_i \in \{32^2, 16^2, 8^2, 4^2\}$ and $C_i \in \{C, 2C, 4C, 8C\}$ for a base width $C$.

The architecture follows standard ResNet~\citep{he2016deep} design, with GroupNorm (GN)~\citep{wu2018group} used in place of BatchNorm (BN)~\citep{ioffe2015batch}.
All residual blocks are ``basic'' blocks (\ie, each consisting of two $3{\times}3$ convolutions).
Following the standard ResNet-34 \citep{he2016deep}: 
the encoder has a 3${\times}$3 convolution (without downsampling) and 4 stages with $[3, 4, 6, 3]$ blocks; downsampling (stride 2) happens at the first block of stages 2 to 4.

For latent-space (\ie, latent-MAE), the input of this ResNet is 32$\times$32$\times$4; for pixel-space, the 256$\times$256$\times$3 input is first patchified (by a 8$\times$8 patch) into 32$\times$32$\times$192. The ResNet operates on the input with $H{\times}W$ $=$ 32$\times$32.

\paragraph{MAE Decoder.}
The decoder returns to the input shape via deconvolutions and skip connections:
\[
\text{Decoder}: \{\mathbf{f}_4, \mathbf{f}_3, \mathbf{f}_2, \mathbf{f}_1\} \mapsto \hat{\mathbf{x}}.
\]

It starts with a $3{\times}3$ convolutional block on $\mathbf{f}_4$, followed by 4 upsampling blocks.
Each upsampling block performs: bilinear 2${\times}$2 upsampling $\to$ concatenating with encoder's skip connection $\to$ GN $\to$ two $3{\times}3$ convolutions with GN and ReLU.
A final 1${\times}$1 convolution produces the output channels. For the pixel-space, the decoder unpatchifies back to the original resolution after the last layer.

\paragraph{Masking.}
The MAE is trained to reconstruct randomly masked inputs.
Unlike the ViT-based MAE~\citep{he2022masked}, which removes the masked tokens from the sequence, we simply zero out masked patches.
For the input of a shape $H{\times}W$ $=$ 32$\times$32 (in either the latent- or pixel-based case), we mask 2${\times}$2 patches by zeroing.
Each patch is independently masked with 50\% probability.

\paragraph{MAE training.}
We minimize the $\ell_2$ reconstruction loss on the masked regions.
We use AdamW \cite{Loshchilov2019} with learning rate $4{\times}10^{-3}$ and a batch size of 8192. EMA with decay 0.9995 is used.
Following \cite{he2022masked}, we apply random resized crop augmentation to the input (for the latent setting, images are augmented before being passed through the VAE encoder).

\paragraph{Classification fine-tuning.}
For our best feature encoder (last row of Table~\ref{tab:ablation_drift_space}), we fine-tune the MAE model with a linear classifier head.
The loss is $\lambda \mathcal{L}_{\text{cls}} + (1-\lambda)\mathcal{L}_{\text{recon}}$. We fine-tune all parameters in this MAE for 3k iterations, where $\lambda$ follows a linear warmup schedule, increasing from $0$ to $0.1$ over the first 1k iterations and remaining constant at $0.1$ for the rest of the training.


%============================================================
% Other Feature Encoders
%============================================================
\subsection{Other Pretrained Feature Encoders}
\label{sec:appendix_other_encoders}

In addition to our customized MAE, we also evaluate other feature encoders for computing the drifting loss.

\paragraph{MoCo and SimCLR.} We evaluate publicly available self-supervised encoders trained on ImageNet in pixel space: MoCo \cite{he2020momentum,chen2020improved} SimCLR \cite{chen2020simple}. We use the ResNet-50 variant. 
For latent-space generation, we apply the VAE decoder to map generator outputs from latent space (32${\times}$32${\times}$4) to pixel space (256${\times}$256${\times}$3) before feature extraction.
Gradients are backpropagated through both the feature extractor and the VAE decoder.

\paragraph{MAE with ConvNeXt-V2.} In our pixel-space generator, we also investigate ConvNeXt-V2 \citep{woo2023convnext} as the feature encoder. We note that ConvNeXt-V2 is a self-supervised pre-trained model using the MAE objective, followed by classification fine-tuning. Like ResNet, ConvNeXt-V2 is a multi-stage architecture.

%============================================================
% Feature Extraction Details
%============================================================
\subsection{Multi-scale Features for Drifting Loss}
\label{sec:appendix_feature_extraction}

Given an image, the feature encoder produces feature maps at multiple scales, with multiple spatial locations per scale.
We compute one drifting loss per feature (\eg, per scale and/or per location).
Specifically, we compute the kernel, the drift, and the resulting loss independently for each feature. The resulting losses are summed.

For each stage in a ResNet, we extract features from the output of every 2 residual blocks, together with the final output.
This yields a set of feature maps, each of shape $H_i{\times}W_i{\times}C_i$. For each feature map, we produce:
\begin{enumerate}[itemsep=2pt, topsep=2pt, parsep=4pt, partopsep=4pt,label=(\alph*)]
    \item $H_i{\times}W_i$ vectors, one per location (each $C_i$-dim);
    \item 1 global mean and 1 global std (each $C_i$-dim);
    \item $\frac{H_i}{2}{\times}\frac{W_i}{2}$ vectors of means and $\frac{H_i}{2}{\times}\frac{W_i}{2}$ vectors of stds (each $C_i$-dim), computed over 2${\times}$2 patches;
    \item $\frac{H_i}{4}{\times}\frac{W_i}{4}$ vectors of means and $\frac{H_i}{4}{\times}\frac{W_i}{4}$ vectors of stds (each $C_i$-dim), computed over 4${\times}$4 patches.
\end{enumerate}
In addition, for the encoder's input ($H_0{\times}W_0{\times}C_0$), we compute the mean of squared values ($x^2$) per channel and obtain a $C_0$-dim vector.

All resulting vectors here are $C_i$-dim.
We compute one drifting loss for each of these $C_i$-dim vectors. All these losses, in addition to the vanilla drifting loss without $\phi$, are summed.
This table compares the effect of these designs on our ablation default:
\begin{center}
\tablestyle{6pt}{1.02}
\tablefontsize
\begin{tabular}{c|ccc}
\toprule
 & (a,b) & (a-c) & (a-d) \\
\midrule
FID & 9.58 & 9.10 & \textbf{8.46} \\
\bottomrule
\end{tabular}
\end{center}
This shows that our method benefits from richer feature sets.
We note that once the feature encoder is run, the computational cost of our drifting loss is negligible: computing multi-scale, multi-location losses incurs little overhead compared to computing a single loss.


\subsection{Feature and Drift Normalization}
\label{sec:appendix_feature_and_loss_normalization}

To balance the multiple loss terms from multiple features, we perform normalization for each feature $\phi_j$, where, $\phi_j$ denotes a feature at a specific spatial location within a given scale (see \ref{sec:appendix_feature_extraction}). 
Intuitively, we want to perform normalization such that the kernel $k(\cdot, \cdot)$ and the drift $\V$ are insensitive to the absolute magnitude of features.
This allows our model to robustly support different feature encoders (see Table~\ref{tab:ablation_drift_space}) as well as a rich set of features from one encoder.

\paragraph{Feature Normalization.} Consider a feature $\phi_j \in \mathbb{R}^{C_j}$. We define a normalization scale $S_j \in \mathbb{R}^{1}$ and the normalized feature is denoted by:
\begin{equation}
\tilde{\phi}_j := {\phi}_j / S_j.
\end{equation}
When using $\tilde{\phi}_j$, the $\ell_2$ distance computed in Eq.~(\ref{eq:kernel}) is:
\begin{equation}
dist_j(\x, \y) = \|\tilde{\phi}_j(\x) - \tilde{\phi}_j(\y)\|,
\end{equation}
where $\x$ denotes a generated sample and $\y$ denotes a positive/negative sample, and $\tilde{\phi}_j(\cdot)$ means extracting their feature at $j$. We want the \textit{average} distance to be $\sqrt{C_j}$:
\begin{equation}
\mathrm{E}_\x \mathrm{E}_\y [ dist_j(\x, \y)] \approx \sqrt{C_j}.
\end{equation}
To achieve this, we set the normalization scale $S_j$ as:
\begin{equation}
S_j = \frac{1}{\sqrt{C_j}} \mathrm{E}_\x \mathrm{E}_\y [ \|{\phi}_j(\x) - {\phi}_j(\y)\| ]
\end{equation}
In practice, we use all $\x$ and $\y$ samples in a batch to compute the empirical mean in place of the expectation. We reuse the \texttt{cdist} computation in Alg.~\ref{alg:drift_code} for computing the pairwise distances. We apply stop-gradient to $S_j$, because this scalar is conceptually computed from samples from the previous batch.

With the normalized feature, the kernel in Eq.~(\ref{eq:kernel}) is set as:
\begin{equation}
k(\x, \y) = \exp\left(-\frac{1}{\tilde{\tau_j}} \|\tilde{\phi}_j(\x) - \tilde{\phi}_j(\y)\|\right),
\end{equation}
where $\tilde{\tau_j}:=\tau{\cdot}\sqrt{C_j}$. By doing so, the value of temperature $\tau$ does not depend on the feature magnitude or feature dimensionality. We set $\tau \in$ \{0.02, 0.05, 0.2\} (discussed next).

\paragraph{Drift Normalization.} When using the feature $\phi_j$, the resulting drift is in the same feature space as $\phi_j$, denoted as $\V_j$. We perform a drift normalization on $\V_j$, for each feature $\phi_j$. Formally, we define a normalization scale $\lambda_j \in \mathbb{R}^1$ and denote:
\begin{equation}
\tilde{\V}_j:=\V_j / \lambda_j.
\end{equation}
Again, we want the normalized drift to be insensitive to the feature magnitude:
\begin{equation}
\mathbb{E} \left[ \frac{1}{C_j}\| \tilde{\V}_j \|^2 \right] \approx 1.
\end{equation}
To achieve this, we set $\lambda_j$ as:
\begin{equation}
\label{eq:lambda}
\lambda_j = \sqrt{\mathbb{E} \left[ \frac{1}{C_j} \| \V_j \|^2 \right]}.
\end{equation}
In practice, the expectation is replaced with the empirical mean computed over the entire batch.

With the normalized feature and normalized drift, the drifting loss of the feature $\phi_j$ is:
\begin{equation}
\label{eq:mse_loss_j}
\mathcal{L}_j = \text{MSE}(\tilde\phi_j({\x}) -  \texttt{sg}(\tilde\phi_j({\x}) + \tilde{\V}_j)),
\end{equation}
where MSE denotes mean squared error.
The overall loss is the sum across all features:
$\mathcal{L} = \sum_j {\mathcal{L}_j}$.

\paragraph{Multiple temperatures.}
Using normalized feature distances, the value of temperature $\tau$ determines \textit{what is considered ``nearby''}.
To improve robustness across different features and across different pretrained models we study, we adopt multiple temperatures.

Formally, for each $\tau$ value, we compute the normalized drift as described above, denoted by $\tilde{\V}_{j,\tau}$. Then we compute an aggregated field: $\tilde{\V}_j \leftarrow \sum_\tau \tilde{\V}_{j,\tau}$, and use it for the loss in \cref{eq:mse_loss_j}. 

% and renormalize it in the same way as Eq.~(\ref{eq:lambda}):
% $\hat{\V}_j \leftarrow \hat{\V}_j / \hat{\lambda}_j$.
% This $\hat{\V}_j$ is used in place of $\tilde{\V}_j$ in ().

This table shows the effect of multiple temperatures on our ablation default:
\begin{center}
    \tablestyle{4pt}{1.05}
    \tablefontsize
    \begin{tabular}{l | x{32} x{32} x{32} | c}
    \toprule
    $\tau$ & 0.02 & 0.05 & 0.2 & \{0.02, 0.05, 0.2\} \\
    \midrule
    FID & 10.62 & \textbf{8.67} & 8.96 & \textbf{8.46} \\
    \bottomrule
    \end{tabular}
\end{center}
Using multiple temperatures can achieve slightly better results than using a single optimal temperature. We fix $\tau \in $ \{0.02, 0.05, 0.2\} and do not require tuning this hyperparameter across different configurations.


\paragraph{Normalization across spatial locations.}
For a feature map of resolution $H_i{\times}W_i$, there are $H_i{\times}W_i$ per-location features.
Separately computing the normalization for each location would be slow and unnecessary.
We assume that features at different locations within the same feature map share the same normalization scale. Accordingly, we concatenate all $H_i{\times}W_i$ locations and compute the normalization scale over all of them.
The feature normalization and drift normalization are both performed in this way.


\subsection{Classifier-Free Guidance (CFG)}
\label{sec:appendix_cfg}

To support CFG, at training time, we include $N_\text{unc}$ additional \mbox{unconditional} samples (real images from random classes) as extra negatives.
These samples are weighted by a factor $w$ when computing the kernel.
For a generated sample $\x$, the effective negative distribution it compares with is:
\[
\tilde{q}(\cdot|c) \triangleq \frac{(N_\text{neg}{-}1) \cdot q_\theta(\cdot|c) + N_\text{unc}  w \cdot p_{\text{data}}(\cdot|\varnothing)}{(N_\text{neg}{-}1) + N_\text{unc} w}.
\]
Comparing this equation with Eq.~(\ref{eq:cfg1})(\ref{eq:cfg2}), we have:
\[
\gamma = \frac{N_\text{unc} w}{(N_\text{neg}{-}1) + N_\text{unc} w}
\]
and 
\[
\alpha = \frac{1}{1-\gamma} = \frac{(N_\text{neg}{-}1) + N_\text{unc} w}{N_\text{neg}{-}1}.
\]
Given a CFG strength $\alpha$, we compute $w$ accordingly, which is used to weight the kernel.
The same weighting $w$ is also applied when computing the global distance normalization.

We train our model with CFG-conditioning \cite{geng2025improved}. At each iteration, we randomly sample $\alpha$ following a pre-defined distribution (see Table~\ref{tab:hypers_all}) and compute the resulting $w$ for weighting the unconditional samples. The value of $\alpha$ is a condition input to the network $f_\theta(\e, c, \alpha)$, alongside the class label $c$.

At inference time, we specify a value of $\alpha$. The inference-time computation remains to be one-step (1-NFE).

\subsection{Sample Queue}
\label{sec:appendix_sample_queue}

Our method requires access to randomly sampled \textit{real} (positive/unconditional) data. This can be implemented using a specialized data loader. 
Instead, we adopt a \textit{sample queue} of cached data, similar to the queue used in MoCo \cite{he2020momentum}.
This implementation samples data in a statistically similar way to a specialized data loader.
For completeness, we describe our implementation as follows, while noting that a data loader would be a more principled solution.  

For each class label, we keep a queue of size 128; for unconditional samples (used in CFG), we maintain a separate global queue of size 1000.
At each training step, we push the latest 64 new real (positive/unconditional) samples, alongside their labels, into the corresponding queues; the earliest ones are dequeued.
When sampling, positive samples are drawn from the queue of the corresponding class, and unconditional samples are drawn from the global queue.
We sample without replacement.

\subsection{Training Loop}
\label{sec:appendix_training_loop}

In summary, in the training loop, each step proceeds as:
\begin{enumerate}[itemsep=2pt, topsep=2pt, parsep=2pt, partopsep=2pt]
    \item Sample a batch ($N_c$) of class labels.
    \item For each label $c$, sample a CFG scale $\alpha$.
    \item Sample a batch ($N_\text{neg}$) of noise $\e$. Feed $(\e, c, \alpha)$ to the generator $f$ to produce generated samples;
    \item Sample positive samples (same class, $N_\text{pos}$) and unconditional samples (for CFG, $N_\text{unc}$);
    \item Extract features on all generated, positive, and unconditional samples
    \item Compute the drifting loss using the features.
    \item Run backpropagation and parameter update.
\end{enumerate}

